\documentclass[12pt,a4paper]{letter}

\usepackage[utf8]{inputenc}
\usepackage[T1]{fontenc}
\usepackage{newtxtext}
\usepackage[margin=2.5cm]{geometry}
\usepackage{hyperref}
\usepackage{xcolor}

\hypersetup{
    colorlinks=true,
    linkcolor=blue,
    urlcolor=blue
}

\signature{Tao Zhu, M.D.}
\address{
    Department of Obstetrics and Gynecology\\
    National Clinical Research Center for Obstetrics and Gynecology\\
    Tongji Hospital, Tongji Medical College\\
    Huazhong University of Science and Technology\\
    Wuhan 430030, China\\
    \medskip
    E-mail: zhutao@tjh.tjmu.edu.cn
}

\begin{document}

\begin{letter}{The Editor-in-Chief\\
\textit{Briefings in Bioinformatics}\\
Oxford University Press}

\opening{Dear Editor,}

I am submitting our manuscript, ``\textbf{Complementary Information Channels for Synthetic Lethality Prediction: Augmenting Graph Topology with Protein Language Models and LLM Weak Supervision},'' for your consideration.

This paper grew out of a straightforward question: can we squeeze more signal out of the biomedical literature than what curated databases already give us? Synthetic lethality prediction has a cold-start problem---all the top methods stumble when asked to predict interactions for genes they have not seen during training. We thought that feeding a GNN both what the literature says about genes (via LLM-driven NER) and what protein sequences look like (via ESM-2 embeddings) might help fill that gap. It did, though not in the way we initially expected.

\medskip
\noindent\textbf{What we found}

The graph topology---PPI edges from STRING---turns out to be the workhorse. SLMGAE, which exploits only graph structure, already hits AUC=0.77 in the hardest cold-start setting. Our added channels (ESM-2 at +12.2\%, LLM text at +1.6\%) are modest supplements, not replacements. We report this honestly in the paper rather than cherry-picking comparisons that make us look better.

But along the way we stumbled into two things that we think are more useful than another incremental performance boost:

\begin{enumerate}
    \item Multi-view GAE methods, including SLMGAE as it is commonly run, leak test-set SL pairs into their training adjacency matrix during cold-start evaluation---because the standard pipeline builds the SL graph from the full database before splitting. When we masked those pairs properly, CV3 AUC dropped from 0.907 to 0.770. That 13.7\% gap is not a model failure; it is a measurement artifact, and we have not seen it discussed in the benchmark literature.

    \item Zero-vector ablation---dropping out a feature channel by replacing it with zeros---is standard practice in GNN papers. It also overestimates feature importance by 14$\times$ in our model, because batch normalization layers do not handle an input of all zeros gracefully. Shuffled-feature ablation gave us a much more honest estimate. We recommend it as a replacement.
\end{enumerate}

\medskip
\noindent\textbf{Why \textit{Briefings in Bioinformatics}}

The paper sits at the intersection of three fast-moving areas---graph neural networks, protein language models, and LLM-driven literature mining---and the methods we used are general enough that researchers working in any of these areas should be able to adapt them. The two measurement issues we flagged are also general: any paper that uses multi-view graph autoencoders for cold-start prediction, or zero-vector ablation in batch-normalized GNNs, stands to benefit from being aware of them. This feels like the kind of broadly useful methodological contribution that \textit{Briefings in Bioinformatics} is known for.

\medskip
\noindent\textbf{Additional details}

The manuscript is single-author. I have no competing interests to declare. All data come from public repositories (PubMed, STRING, SynLethDB, DepMap). The code and processed data are at \url{https://github.com/tjogzt/LLM4SL} (DOI: 10.5281/zenodo.21743293). This work has not been submitted elsewhere.

\medskip
Thank you for your time and consideration.

\medskip
\closing{Sincerely,}

\end{letter}

\end{document}
